% 2006--2007 reviewed draft.
% Source PDFs:
% - elegantbook2/试卷/2006数学分析下A卷.pdf
% - elegantbook2/试卷/2006数学分析下B卷.pdf
% - elegantbook2/试卷/2007数学分析下A卷.pdf
% - elegantbook2/试卷/2007数学分析下A卷及答案.pdf
% Raw OCR/text sources under word/past-exams/raw were used as auxiliary references;
% low-confidence OCR points have been resolved against page images where needed.

\section{2006--2007 第二学期 A 卷}

\subsection{试题}

\paragraph{一、单项选择题（每小题 3 分，共 15 分）}

\begin{enumerate}
  \item 设有直线
  \[
    L:\begin{cases}
      x+3y+2z+1=0,\\
      2x-y-10z+3=0
    \end{cases}
  \]
  及平面 $\pi:4x-2y+z-2=0$，则直线 $L$（\quad）。

  A. 平行于平面 $\pi$；\quad
  B. 在平面 $\pi$ 上；\quad
  C. 垂直于平面 $\pi$；\quad
  D. 与平面 $\pi$ 斜交。

  \item 二元函数
  \[
    f(x,y)=
    \begin{cases}
      \dfrac{x^2y}{x^2+y^2}, & (x,y)\ne(0,0),\\
      0, & (x,y)=(0,0)
    \end{cases}
  \]
  在点 $(0,0)$ 处（\quad）。

  A. 偏导数存在，但不可微；\quad
  B. 连续、偏导数不存在；\quad
  C. 不连续、偏导数不存在；\quad
  D. 偏导数不存在。

  \item 设平面区域
  \[
    D=\{(x,y)\mid -a\le x\le a,\ x\le y\le a\},\qquad
    D_1=\{(x,y)\mid 0\le x\le a,\ x\le y\le a\},
  \]
  则
  \[
    \iint_D (xy+\cos x\sin y)\,dxdy
  \]
  等于（\quad）。

  A. $2\iint_{D_1}\cos x\sin y\,dxdy$；\quad
  B. $2\iint_{D_1}xy\,dxdy$；\quad
  C. $4\iint_{D_1}(xy+\cos x\sin y)\,dxdy$；\quad
  D. $0$。

  \item 若级数 $\sum_{n=1}^{\infty}a_n^2$，$\sum_{n=1}^{\infty}b_n^2$ 都收敛，则下列结论不成立的是（\quad）。

  A. $\sum_{n=1}^{\infty}(a_n+b_n)^2$ 收敛；\quad
  B. $\sum_{n=1}^{\infty}\dfrac{|a_n|}{n}$ 收敛；\quad
  C. $\sum_{n=1}^{\infty}a_nb_n$ 收敛；\quad
  D. $\sum_{n=1}^{\infty}a_nb_n$ 发散。

  \item 微分方程
  \[
    y'+\frac{1}{x}y=\frac{\sin x}{x}
  \]
  的通解为（\quad）。

  A. $y=\dfrac{1}{x}(-\cos x+C)$；\quad
  B. $y=x(-\cos x+C)$；\quad
  C. $y=\dfrac{1}{x}(-\sin x+C)$；\quad
  D. $y=\dfrac{1}{x}(\cos x+C)$。
\end{enumerate}

\paragraph{二、填空题（每小题 3 分，共 15 分）}

\begin{enumerate}
  \item 过点 $(0,2,4)$ 且同时平行于平面 $x+2z=1$ 和 $y-3z=2$ 的直线方程是 \underline{\hspace{4cm}}。

  \item 将积分
  \[
    \int_1^e dx\int_0^{\ln x} f(x,y)\,dy
  \]
  改变积分顺序为 \underline{\hspace{4cm}}。
  \textit{（校勘：原卷内层微分疑似误印为 $dx$；按二重积分换序题意应为 $dy$。）}

  \item 级数
  \[
    \sum_{n=1}^{\infty}\frac{3^n+(-2)^n}{n}(x+1)^n
  \]
  的收敛域为 \underline{\hspace{4cm}}。

  \item 设函数 $u=xyz$，则 $u$ 沿方向 $\ell=\{2,-1,3\}$ 的方向导数为 \underline{\hspace{4cm}}。

  \item 设 $L$ 为取逆时针方向的圆周 $x^2+y^2=9$，求曲线积分
  \[
    \oint_L (2xy-2y)\,dx+(x^2-4x)\,dy
  \]
  的值为 \underline{\hspace{4cm}}。
\end{enumerate}

\paragraph{三、计算题（每小题 7 分，共 35 分）}

\begin{enumerate}
  \item 将函数
  \[
    f(x)=\ln\frac{1}{2+2x+x^2}
  \]
  按 $(x+1)$ 的正整数次幂展开成幂级数，即在 $x_0=-1$ 处展开成泰勒级数。

  \item 利用高斯公式计算曲面积分
  \[
    \iint_{\Sigma}x^2\,dydz+2y^2\,dzdx+(3z^2-4x^2y^2)\,dxdy,
  \]
  其中 $\Sigma$ 为 $z=\sqrt{x^2+y^2}$ 与 $z=2$ 围成的立体表面，取外侧。

  \item 设二阶可微函数 $\varphi(x)$ 满足方程
  \[
    \varphi(x)=e-\int_0^x (x-u)\varphi(u)\,du,
  \]
  求 $\varphi(x)$。

  \item 设
  \[
    z=xf\left(2x,\frac{y^2}{x}\right),
  \]
  其中 $f$ 具有二阶连续偏导数，求 $\dfrac{\partial^2 z}{\partial x\partial y}$。

  \item 设三重积分
  \[
    I=\iiint_V (x^2+y^2+z^2)\,dxdydz,
  \]
  其中积分区域 $V$ 是由球面 $x^2+y^2+z^2=R^2$（$z\ge 0$）及锥面 $x^2+y^2=z^2$ 所围成的区域。
\end{enumerate}

\paragraph{四、证明题（第 1 题 7 分，第 2 题 8 分，共 15 分）}

\begin{enumerate}
  \item 证明：一般项级数
  \[
    \sum_{n=1}^{\infty}\sin\left(\pi\sqrt{n^2+1}\right)
  \]
  条件收敛。

  \item 证明：曲面
  \[
    F\left(\frac{x-a}{z-c},\frac{y-b}{z-c}\right)=0
  \]
  上任意一点的切平面都通过一定点。
\end{enumerate}

\paragraph{五、应用题（每小题 10 分，共 20 分）}

\begin{enumerate}
  \item 设函数 $y(x)$（$x\ge 0$）二阶可导，且 $y'(x)>0$，$y(0)=1$。
  过曲线 $y=y(x)$ 上任意一点 $P(x,y)$ 作该曲线的切线及 $x$ 轴的垂线，上述两直线与 $x$ 轴所围成的三角形面积记为 $S_1$。
  区间 $(0,x]$ 上以 $y=y(x)$ 为曲边的曲边梯形面积记为 $S_2$。
  若 $2S_1-S_2$ 恒为 $1$，求此曲线 $y=y(x)$ 的方程。

  \item 将长为 $a$ 的线段截成两段，一段围成一个正方形，另一段围成一个圆。问这两段的长度分别为多少时，它们所围成的正方形和圆的面积之和最大？
\end{enumerate}

\subsection{答案与解析}

\paragraph{一、单项选择题}
\begin{enumerate}
  \item 由两平面法向量 $(1,3,2)$ 与 $(2,-1,-10)$ 叉乘，得直线方向向量可取 $(4,-2,1)$，它与平面 $\pi$ 的法向量 $(4,-2,1)$ 平行，故直线垂直于平面。选 C。
  \item 函数在原点连续，且两个偏导数均存在；但沿 $y=x$ 时
  \[
  \frac{f(x,x)-f(0,0)}{\sqrt{x^2+x^2}}=\frac{x/2}{\sqrt2|x|}
  \]
  不趋于 $0$，故不可微。选 A。
  \item 直接积分可得
  \[
  \iint_Dxy\,dxdy=0,\qquad
  \iint_D\cos x\sin y\,dxdy=2\iint_{D_1}\cos x\sin y\,dxdy.
  \]
  选 A。
  \item A、B、C 均可由 Cauchy 不等式或比较判别推出收敛，D 与 C 矛盾，故不成立。选 D。
  \item 方程的积分因子为 $x$，于是
  \[
  (xy)'=\sin x,\qquad xy=-\cos x+C.
  \]
  选 A。
\end{enumerate}

\paragraph{二、填空题}
\begin{enumerate}
  \item 两平面的法向量为 $(1,0,2)$ 与 $(0,1,-3)$，直线方向向量可取
  \[
  (1,0,2)\times(0,1,-3)=(-2,3,1).
  \]
  因此直线方程为
  \[
  \frac{x}{-2}=\frac{y-2}{3}=z-4.
  \]
  \item 积分区域为 $1\le x\le e,\ 0\le y\le\ln x$，即
  \[
  0\le y\le1,\qquad e^y\le x\le e.
  \]
  故
  \[
  \int_1^e dx\int_0^{\ln x}f(x,y)\,dy
  =\int_0^1dy\int_{e^y}^e f(x,y)\,dx.
  \]
  \item 半径由 $3^n$ 项决定，为 $1/3$，中心为 $-1$。当 $x+1=-1/3$ 时为交错调和级数加收敛级数；当 $x+1=1/3$ 时含调和级数发散。收敛域为
  \[
  \left[-\frac43,-\frac23\right).
  \]
  \item
  \[
  \nabla u=(yz,xz,xy),\qquad
  \frac{\ell}{|\ell|}=\frac{(2,-1,3)}{\sqrt{14}},
  \]
  故方向导数为
  \[
  D_\ell u=\frac{2yz-xz+3xy}{\sqrt{14}}.
  \]
  \item 由 Green 公式，
  \[
  Q_x-P_y=(2x-4)-(2x-2)=-2.
  \]
  圆盘面积为 $9\pi$，故积分为
  \[
  -18\pi.
  \]
\end{enumerate}

\paragraph{三、计算题}
\begin{enumerate}
  \item 令 $t=x+1$，则
  \[
  f(x)=\ln\frac1{1+t^2}=-\ln(1+t^2)
  =\sum_{n=1}^{\infty}\frac{(-1)^n}{n}t^{2n}.
  \]
  因此
  \[
  f(x)=\sum_{n=1}^{\infty}\frac{(-1)^n}{n}(x+1)^{2n},\qquad |x+1|<1.
  \]
  \item 设向量场
  \[
  \mathbf F=(x^2,2y^2,3z^2-4x^2y^2).
  \]
  其散度为
  \[
  \nabla\cdot\mathbf F=2x+4y+6z.
  \]
  圆锥体关于 $x,y$ 对称，$2x,4y$ 的积分为 $0$。在柱坐标下区域为 $0\le z\le2,\ 0\le r\le z$，故
  \[
  \iint_{\Sigma}\mathbf F\cdot d\mathbf S
  =\int_0^{2\pi}\int_0^2\int_0^z6z\,r\,dr\,dz\,d\theta
  =24\pi.
  \]
  \item 对方程
  \[
  \varphi(x)=e-\int_0^x(x-u)\varphi(u)\,du
  \]
  求导，得 $\varphi'(x)=-\int_0^x\varphi(u)\,du$，再求导得
  \[
  \varphi''(x)=-\varphi(x),\qquad \varphi(0)=e,\quad \varphi'(0)=0.
  \]
  因此
  \[
  \varphi(x)=e\cos x.
  \]
  \item 设 $u=2x,\ v=y^2/x$。先对 $y$ 求导：
  \[
  z_y=x f_2(u,v)\frac{2y}{x}=2y f_2(u,v).
  \]
  再对 $x$ 求导：
  \[
  z_{xy}=2y\left(2f_{21}(u,v)-\frac{y^2}{x^2}f_{22}(u,v)\right).
  \]
  即
  \[
  \frac{\partial^2z}{\partial x\partial y}
  =4y f_{12}\left(2x,\frac{y^2}{x}\right)
  -\frac{2y^3}{x^2}f_{22}\left(2x,\frac{y^2}{x}\right).
  \]
  \item 用球坐标。锥面 $x^2+y^2=z^2$ 对应 $\varphi=\pi/4$，区域为
  \[
  0\le\rho\le R,\qquad 0\le\varphi\le\frac{\pi}{4},\qquad 0\le\theta\le2\pi.
  \]
  因 $x^2+y^2+z^2=\rho^2$，故
  \[
  I=\int_0^{2\pi}\int_0^{\pi/4}\int_0^R\rho^4\sin\varphi\,d\rho d\varphi d\theta
  =\frac{\pi R^5}{5}(2-\sqrt2).
  \]
\end{enumerate}

\paragraph{四、证明题}
\begin{enumerate}
  \item 记
  \[
  \delta_n=\pi\left(\sqrt{n^2+1}-n\right)\sim\frac{\pi}{2n}.
  \]
  则
  \[
  \sin\left(\pi\sqrt{n^2+1}\right)
  =(-1)^n\sin\delta_n.
  \]
  由交错级数判别法，原级数收敛；又
  \[
  \left|\sin\left(\pi\sqrt{n^2+1}\right)\right|\sim\frac{\pi}{2n},
  \]
  绝对值级数发散，故原级数条件收敛。
  \item 设
  \[
  u=\frac{x-a}{z-c},\qquad v=\frac{y-b}{z-c}.
  \]
  曲面由过定点 $(a,b,c)$ 的直线族
  \[
  (x,y,z)=(a,b,c)+t(u,v,1)
  \]
  生成。经过曲面上一点并沿该母线的方向仍在曲面上，所以该母线方向属于该点切平面，切平面必通过定点 $(a,b,c)$。
\end{enumerate}

\paragraph{五、应用题}
\begin{enumerate}
  \item 切线与 $x$ 轴交点到 $x$ 的水平距离为 $y/y'$，故
  \[
  S_1=\frac12y\cdot\frac{y}{y'}=\frac{y^2}{2y'},\qquad
  S_2=\int_0^x y(t)\,dt.
  \]
  条件 $2S_1-S_2=1$ 化为
  \[
  \frac{y^2}{y'}-\int_0^xy(t)\,dt=1.
  \]
  求导得
  \[
  yy''=(y')^2.
  \]
  因此 $(y'/y)'=0$，即 $y=Ce^{kx}$。由 $y(0)=1$ 得 $C=1$；再代入 $x=0$ 的原条件得 $1/y'(0)=1$，故 $k=1$。所以
  \[
  y=e^x.
  \]
  \item 设围正方形的线段长为 $s$，围圆的线段长为 $a-s$。面积和为
  \[
  A(s)=\frac{s^2}{16}+\frac{(a-s)^2}{4\pi},\qquad 0\le s\le a.
  \]
  这是开口向上的二次函数，最大值在端点取得。比较
  \[
  A(0)=\frac{a^2}{4\pi},\qquad A(a)=\frac{a^2}{16},
  \]
  因 $1/(4\pi)>1/16$，故最大时正方形段为 $0$，圆段为 $a$。若要求两段都为正，则没有最大值，只有上确界 $a^2/(4\pi)$。
\end{enumerate}

\section{2006--2007 第二学期 B 卷}

\subsection{试题}

\paragraph{一、单项选择题（每小题 3 分，共 15 分）}

\begin{enumerate}
  \item 下列说法不正确的是（\quad）。

  A. 如果函数 $f(x,y)$ 在区域 $D$ 中连续，则 $|f(x,y)|$ 也在 $D$ 中连续；\quad
  B. 函数 $f(x,y)$ 沿射线 $\ell$ 的方向导数，等于梯度在 $\ell$ 上的投影；\quad
  C. 如果函数 $f(x,y)$ 在点 $P(x_0,y_0)$ 取极值，则必有 $f_x(x_0,y_0)=0,\ f_y(x_0,y_0)=0$；\quad
  D. 如果函数 $f(x,y)$ 在点 $P(x_0,y_0)$ 存在两个连续的偏导数，则在该点一定可微。

  \item 函数 $z=\ln(1+x^2+y^2)$ 当 $x=1,\ y=2$ 时的全微分为（\quad）。

  A. $\dfrac13dx+\dfrac13dy$；\quad
  B. $\dfrac23dx+\dfrac13dy$；\quad
  C. $\dfrac23dx+\dfrac23dy$；\quad
  D. $\dfrac13dx+\dfrac23dy$。

  \item 设 $\Omega$ 是曲面 $z=x^2+y^2$ 与平面 $z=1$ 所围成的闭区域，则 $\iiint_{\Omega}f(x,y,z)\,dV$ 化为三次积分为（\quad）。

  A. $\displaystyle\int_{-1}^{1}dy\int_{-\sqrt{1-y^2}}^{\sqrt{1-y^2}}dx\int_1^{x^2+y^2}f(x,y,z)\,dz$；

  B. $\displaystyle4\int_0^1dx\int_0^{\sqrt{1-x^2}}dy\int_{x^2+y^2}^{1}f(x,y,z)\,dz$；

  C. $\displaystyle4\int_0^1dx\int_0^{\sqrt{1-x^2}}dy\int_0^{x^2+y^2}f(x,y,z)\,dz$；

  D. $\displaystyle\int_{-1}^{1}dx\int_{-\sqrt{1-x^2}}^{\sqrt{1-x^2}}dy\int_{x^2+y^2}^{1}f(x,y,z)\,dz$。

  \item 设曲线 $L$ 为沿顺时针方向的圆周 $x^2+y^2=9$，则第二型闭路曲线积分
  \[
    \oint_L (2xy-2y)\,dx+(x^2-4x)\,dy
  \]
  等于（\quad）。

  A. $0$；\quad B. $2\pi$；\quad C. $6\pi$；\quad D. $18\pi$。

  \item 已知幂级数 $\sum_{n=1}^{\infty}a_nx^n$ 的收敛域为 $[-8,8]$，则
  \[
    \sum_{n=2}^{\infty}\frac{a_nx^n}{n(n-1)}
  \]
  的收敛半径为（\quad）。
  \textit{（校勘：原卷下标疑似为 $n=1$，但 $n(n-1)$ 使首项无定义；此处按 $n=2$ 处理。）}

  A. $R>8$；\quad B. $R<8$；\quad C. $R=8$；\quad D. 不能确定。
\end{enumerate}

\paragraph{二、填空题（每小题 3 分，共 15 分）}

\begin{enumerate}
  \item 点 $(2,1,0)$ 到平面 $3x+4y+5z=0$ 的距离为 \underline{\hspace{4cm}}。

  \item 微分方程
  \[
    \frac{1}{y}\frac{dy}{dx}=2x+\frac{x-x^3}{y}
  \]
  （$y\ne0$）的通解为 \underline{\hspace{4cm}}。

  \item 设
  \[
    f(x,y)=\ln\left(x+\frac{y}{2x}\right),
  \]
  则 $\left.\dfrac{\partial f}{\partial y}\right|_{x=1,\ y=0}=$ \underline{\hspace{4cm}}。

  \item 设
  \[
    f(x)=
    \begin{cases}
      x^2, & -2\le x<0,\\
      4-x^2, & 0\le x<2
    \end{cases}
  \]
  则 $f(x)$ 以 $4$ 为周期的傅里叶级数的和函数 $s(x)$ 在 $[-2,2)$ 上的表达式为 \underline{\hspace{4cm}}。

  \item 设曲线积分
  \[
    \int_L [f(x)-e^x]\sin y\,dx-f(x)\cos y\,dy
  \]
  与积分路径无关，其中 $f(x)$ 具有一阶连续导数，且 $f(0)=0$，则 $f(x)=$ \underline{\hspace{4cm}}。
\end{enumerate}

\paragraph{三、计算题（每小题 7 分，共 35 分）}

\begin{enumerate}
  \item 判断交错级数
  \[
    \sum_{n=2}^{\infty}\frac{(-1)^n}{\sqrt n+(-1)^n}
  \]
  的敛散性。
  \textit{（校勘：原卷下标疑似为 $n=1$，但首项分母为 $0$；此处按 $n=2$ 处理。）}

  \item 利用高斯公式计算曲面积分
  \[
    \iint_{\Sigma}x^2\,dydz+2y^2\,dzdx+(3z^2-4x^2y^2)\,dxdy,
  \]
  其中 $\Sigma$ 为 $z=\sqrt{x^2+y^2}$ 与 $z=2$ 围成的立体表面，取外侧。

  \item 已知 $y_1(x)=e^x$ 是齐次线性微分方程
  \[
    (2x-1)y''-(2x+1)y'+2y=0
  \]
  的一个解，求此方程的通解。

  \item 设
  \[
    z=xf\left(\frac{y}{x}\right)+yg\left(x,\frac{x}{y}\right),
  \]
  其中 $f,g$ 均为二阶可微函数，求 $\dfrac{\partial^2 z}{\partial x\partial y}$。

  \item 计算二重积分
  \[
    \iint_D (x^2+y^2)\,d\sigma,
  \]
  其中 $D$ 是由 $y=\sqrt{1-x^2}$ 和 $x$ 轴所围成的平面区域。
\end{enumerate}

\paragraph{四、证明题（第 1 题 7 分，第 2 题 8 分，共 15 分）}

\begin{enumerate}
  \item 证明：级数 $\sum_{n=1}^{\infty}(a_{n+1}-a_n)$ 收敛的充分必要条件是数列 $\{a_n\}$ 收敛。

  \item 证明：曲面
  \[
    F\left(\frac{x-a}{z-c},\frac{y-b}{z-c}\right)=0
  \]
  上任意一点的切平面都通过一定点。
\end{enumerate}

\paragraph{五、应用题（每小题 10 分，共 20 分）}

\begin{enumerate}
  \item 设有力场
  \[
    \mathbf F=\left(-\frac1y,\frac{x}{y^2}\right),
  \]
  证明力 $\mathbf F$ 在上半平面内对质点所做的功与路径无关，并计算质点从点 $A(1,2)$ 移动到点 $B(2,1)$ 时力 $\mathbf F$ 所做的功。

  \item 设周长为 $2p$ 的矩形绕它的一边旋转构成圆柱形，求矩形的边长各为多少时，圆柱的体积最大。
\end{enumerate}

\subsection{答案与解析}

\paragraph{一、单项选择题}
\begin{enumerate}
  \item C 不正确。极值点若偏导不存在，不能推出 $f_x=f_y=0$。选 C。
  \item
  \[
  dz=\frac{2x}{1+x^2+y^2}dx+\frac{2y}{1+x^2+y^2}dy,
  \]
  在 $(1,2)$ 处为 $\dfrac13dx+\dfrac23dy$。选 D。
  \item 区域满足 $x^2+y^2\le1,\ x^2+y^2\le z\le1$，故选 D。
  \item 对顺时针方向，Green 公式取负号。因
  \[
  Q_x-P_y=(2x-4)-(2x-2)=-2,
  \]
  逆时针积分为 $-18\pi$，顺时针积分为 $18\pi$。选 D。
  \item 除以 $n(n-1)$ 不改变幂级数收敛半径，故 $R=8$。选 C。
\end{enumerate}

\paragraph{二、填空题}
\begin{enumerate}
  \item
  \[
  d=\frac{|3\cdot2+4\cdot1+5\cdot0|}{\sqrt{3^2+4^2+5^2}}=\sqrt2.
  \]
  \item 原方程化为
  \[
  y'=2xy+x-x^3.
  \]
  即
  \[
  y'-2xy=x-x^3.
  \]
  取积分因子 $e^{-x^2}$，得
  \[
  (ye^{-x^2})'=(x-x^3)e^{-x^2}.
  \]
  积分后通解为
  \[
  y=\frac{x^2}{2}+Ce^{x^2}.
  \]
  \item
  \[
  f_y(x,y)=\frac{1}{x+y/(2x)}\cdot\frac1{2x},
  \]
  故
  \[
  \left.f_y\right|_{(1,0)}=\frac12.
  \]
  \item 傅里叶级数在连续点收敛到函数值，在跳跃点收敛到左右极限平均值。因此
  \[
  s(x)=
  \begin{cases}
  x^2,&-2<x<0,\\
  4-x^2,&0<x<2,\\
  2,&x=-2,0.
  \end{cases}
  \]
  \item 令
  \[
  P=(f(x)-e^x)\sin y,\qquad Q=-f(x)\cos y.
  \]
  路径无关要求 $P_y=Q_x$，即
  \[
  (f-e^x)\cos y=-f'(x)\cos y.
  \]
  故 $f'+f=e^x$，结合 $f(0)=0$ 得
  \[
  f(x)=\frac{e^x-e^{-x}}2.
  \]
\end{enumerate}

\paragraph{三、计算题}
\begin{enumerate}
  \item 当 $n\ge2$ 时，
  \[
  \frac{(-1)^n}{\sqrt n+(-1)^n}
  =
  \frac{(-1)^n(\sqrt n-(-1)^n)}{n-1}
  =
  \frac{(-1)^n\sqrt n}{n-1}-\frac1{n-1}.
  \]
  第一项为交错型且趋于 $0$，第二项为负调和型，故级数发散到 $-\infty$。
  \item 与 A 卷同题。由 Gauss 公式，
  \[
  \iint_{\Sigma}x^2\,dydz+2y^2\,dzdx+(3z^2-4x^2y^2)\,dxdy=24\pi.
  \]
  \item 已知一解 $y_1=e^x$。令 $y=ve^x$，代入方程可化为
  \[
  (2x-1)v''+(2x-3)v'=0.
  \]
  设 $w=v'$，则
  \[
  (2x-1)w'+(2x-3)w=0,
  \]
  得 $w=C(2x-1)e^{-x}$，积分得另一解可取 $2x+1$。故通解为
  \[
  y=C_1e^x+C_2(2x+1).
  \]
  \item 直接链式求导可得
  \[
  \frac{\partial^2z}{\partial x\partial y}
  =-\frac{y}{x^2}f''\left(\frac yx\right)
  +g_1\left(x,\frac xy\right)
  -\frac{x}{y}g_{12}\left(x,\frac xy\right)
  -\frac{x}{y^2}g_{22}\left(x,\frac xy\right).
  \]
  \item 区域为上半单位圆盘，取极坐标：
  \[
  \iint_D(x^2+y^2)\,d\sigma
  =\int_0^\pi\int_0^1 r^2\cdot r\,dr\,d\theta
  =\frac{\pi}{4}.
  \]
\end{enumerate}

\paragraph{四、证明题}
\begin{enumerate}
  \item 部分和为
  \[
  \sum_{n=1}^N(a_{n+1}-a_n)=a_{N+1}-a_1.
  \]
  因此级数收敛当且仅当 $\{a_{N+1}\}$ 收敛，也即 $\{a_n\}$ 收敛。
  \item 同 A 卷：曲面由过定点 $(a,b,c)$ 的母线族生成，母线方向属于切平面，所以任意切平面都通过 $(a,b,c)$。
\end{enumerate}

\paragraph{五、应用题}
\begin{enumerate}
  \item 设
  \[
  P=-\frac1y,\qquad Q=\frac{x}{y^2}.
  \]
  有 $P_y=Q_x=1/y^2$，上半平面单连通，故路径无关。势函数可取
  \[
  \phi(x,y)=-\frac{x}{y}.
  \]
  因此从 $A(1,2)$ 到 $B(2,1)$ 所做的功为
  \[
  \phi(B)-\phi(A)=-2-\left(-\frac12\right)=-\frac32.
  \]
  \item 设矩形一边为 $x$，另一边为 $p-x$。若绕边 $x$ 旋转，则体积
  \[
  V=\pi(p-x)^2x.
  \]
  求导得最大值在 $x=p/3$ 处取得，此时另一边为 $2p/3$。故矩形边长为
  \[
  \frac p3,\qquad \frac{2p}{3},
  \]
  且绕较短边旋转时体积最大，最大体积为
  \[
  \frac{4\pi p^3}{27}.
  \]
\end{enumerate}

\section{2007--2008 第二学期 A 卷}

\subsection{试题}

\paragraph{一、填空题（共 5 小题，每小题 3 分，共 15 分）}

\begin{enumerate}
  \item 函数 $z=\arctan\dfrac{y}{x}$ 在点 $(1,0)$ 处的梯度为 \underline{\hspace{3cm}}。

  \item 积分
  \[
    \int_0^2 dx\int_x^2 e^{-y^2}\,dy
    =
    \underline{\hspace{3cm}}
  \]
  。

  \item 幂级数
  \[
    \sum_{n=1}^{\infty}\frac{(x-2)^n}{n4^n}
  \]
  的收敛域为 \underline{\hspace{3cm}}。

  \item 设数量场 $u=\ln\sqrt{x^2+y^2}$，则
  \[
    \operatorname{div}(\operatorname{grad}u)=\underline{\hspace{3cm}}
  \]
  。

  \item 函数 $z=xy$ 在条件 $x+y=1$ 下的极大值为 \underline{\hspace{3cm}}。
\end{enumerate}

\paragraph{二、选择题（共 5 小题，每小题 3 分，共 15 分）}

\begin{enumerate}
  \item 抛物线 $y^2=4x$ 上与直线 $x-y+4=0$ 相距最近的点是（\quad）。

  A. $(0,0)$；\quad B. $(1,1)$；\quad C. $(1,2)$；\quad D. $(2,2\sqrt2)$。

  \item 设区域 $D:0\le y\le x,\ 0\le x\le \pi$，则
  \[
    \iint_D \sqrt{1-\cos^2 x}\,dxdy
  \]
  等于（\quad）。

  A. $0$；\quad B. $\dfrac{\pi}{2}$；\quad C. $2$；\quad D. $\pi$。

  \item 设 $\alpha$ 是常数，则级数
  \[
    \sum_{n=1}^{\infty}\left[\frac{\sin(n\alpha)}{n^2}-\frac1{\sqrt n}\right]
  \]
  是（\quad）。

  A. 绝对收敛；\quad B. 发散；\quad C. 条件收敛；\quad D. 敛散性与 $\alpha$ 取值有关。

  \item 若 $0\le a_n<\dfrac1n\ (n=1,2,\cdots)$，则下列级数中肯定收敛的是（\quad）。

  A. $\sum_{n=1}^{\infty}a_n$；\quad
  B. $\sum_{n=1}^{\infty}(-1)^na_n$；\quad
  C. $\sum_{n=1}^{\infty}\sqrt{a_n}$；\quad
  D. $\sum_{n=1}^{\infty}(-1)^na_n^2$。

  \item 设 $y_1$ 与 $y_2$ 是非齐次线性微分方程
  \[
    y''+ay'+by=f(x)
  \]
  的两个特解，则下列说法正确的是（\quad）。

  A. $y_1+y_2$ 是非齐次方程的解；\quad
  B. $y_1-y_2$ 是非齐次方程的解；\quad
  C. $y_1+y_2$ 是对应齐次方程的解；\quad
  D. $y_1-y_2$ 是对应齐次方程的解。
\end{enumerate}

\paragraph{三、计算题（共 2 小题，每小题 7 分，共 14 分）}

\begin{enumerate}
  \item 设 $u(x,y,z)=x+yz$，其中 $z=z(x,y)$ 由 $x+y+z=0$ 确定，求
  \[
    \frac{\partial^2u}{\partial x\partial y}.
  \]

  \item 求曲面 $x^2+2y^2+3z^2=21$ 上平行于平面 $x+4y+6z=0$ 的各个切平面。
\end{enumerate}

\paragraph{四、计算题（共 2 小题，每小题 7 分，共 14 分）}

\begin{enumerate}
  \item 计算曲线积分
  \[
    \int_L xy\,dx+(y-x)\,dy,
  \]
  其中 $L$ 为 $y=x^2$ 上由点 $O(0,0)$ 到点 $A(1,1)$ 的一段。

  \item 设 $\Sigma$ 是锥面 $z=\sqrt{x^2+y^2}$ 被平面 $z=1$ 所截下的有限部分下侧，计算
  \[
    \iint_{\Sigma}x\,dydz+y\,dzdx+(4-2z)\,dxdy.
  \]
\end{enumerate}

\paragraph{五、证明题（8 分）}

将函数 $f(x)=x$ 在 $[0,\pi]$ 上展开成余弦级数，并证明
\[
  \sum_{n=1}^{\infty}\frac1{(2n-1)^2}=\frac{\pi^2}{8}.
\]

\paragraph{六、计算题（8 分）}

求幂级数
\[
  \sum_{n=1}^{\infty}n x^{2n}
\]
的收敛域及和函数。

\paragraph{七、证明题（8 分）}

设 $f(x)\in C[0,2]$，证明
\[
  \int_0^2 dx\int_x^2 f(x)f(y)\,dy
  =
  \frac12\left[\int_0^2 f(x)\,dx\right]^2.
\]

\paragraph{八、证明题（8 分）}

证明函数
\[
  F(y)=\int_1^{+\infty}\frac{dx}{\sqrt[3]{x^4+y^2}}
\]
在 $-\infty<y<+\infty$ 上连续。

\paragraph{九、应用题（10 分）}

设
\[
  f(x)=\sin x-\int_0^x (x-t)f(t)\,dt,
\]
其中 $f(x)$ 为连续函数，求 $f(x)$。

\subsection{答案与解析}

\paragraph{一、填空题}

\begin{enumerate}
  \item $\operatorname{grad}z(1,0)=(0,1)$。
  \item 交换积分次序：
  \[
    \int_0^2dx\int_x^2e^{-y^2}\,dy
    =
    \int_0^2dy\int_0^y e^{-y^2}\,dx
    =
    \int_0^2ye^{-y^2}\,dy
    =
    \frac{1-e^{-4}}{2}
  \]
  。
  \textit{（校勘：答案版给出的数值对应 $e^{y^2}$，但试题 PDF 图像为 $e^{-y^2}$；此处按题面计算。）}
  \item 令 $t=\dfrac{x-2}{4}$，则级数为 $\displaystyle\sum_{n=1}^\infty \frac{t^n}{n}$。当 $t=-1$ 时收敛，当 $t=1$ 时发散，故收敛域为 $[-2,6)$。
  \item $\operatorname{div}(\operatorname{grad}u)=0$（在 $(x,y)\ne(0,0)$ 上）。
  \item 极大值为 $\dfrac14$。
\end{enumerate}

\paragraph{二、选择题}

\begin{enumerate}
  \item 设抛物线上点为 $\left(\dfrac{y^2}{4},y\right)$，到直线的距离与
  \[
    \frac{y^2}{4}-y+4=\frac{(y-2)^2}{4}+3
  \]
  同时取最小。故 $y=2,\ x=1$，选 C。
  \item 在 $0\le x\le\pi$ 上有 $\sqrt{1-\cos^2x}=\sin x$，故
  \[
    \iint_D\sqrt{1-\cos^2x}\,dxdy=\int_0^\pi x\sin x\,dx=\pi.
  \]
  选 D。
  \item 因 $\displaystyle\sum_{n=1}^\infty \left|\frac{\sin(n\alpha)}{n^2}\right|$ 收敛，而 $\displaystyle\sum_{n=1}^\infty\frac1{\sqrt n}$ 发散，所以原级数发散，选 B。
  \item 由 $0\le a_n<\dfrac1n$ 得 $0\le a_n^2<\dfrac1{n^2}$，所以 $\displaystyle\sum(-1)^na_n^2$ 绝对收敛，选 D。
  \item 设线性算子 $L[y]=y''+ay'+by$。因 $L[y_1]=L[y_2]=f(x)$，故 $L[y_1-y_2]=0$，所以 $y_1-y_2$ 是对应齐次方程的解，选 D。
\end{enumerate}

\paragraph{三、计算题}

\begin{enumerate}
  \item 由 $x+y+z=0$ 得 $z_x=z_y=-1$，$z_{xy}=0$。又 $u=x+yz$，故
  \[
    \frac{\partial^2u}{\partial x\partial y}
    =
    \frac{\partial}{\partial y}\left(1+yz_x\right)
    =
    \frac{\partial}{\partial y}(1-y)=-1.
  \]

  \item 设切点为 $(x_0,y_0,z_0)$。椭球面法向量为
  \[
    (2x_0,4y_0,6z_0),
  \]
  与平面 $x+4y+6z=0$ 的法向量 $(1,4,6)$ 平行，因此
  \[
    (2x_0,4y_0,6z_0)=\lambda(1,4,6).
  \]
  代入 $x_0^2+2y_0^2+3z_0^2=21$ 得 $\lambda=\pm2$，故切平面为
  \[
    x+4y+6z=\pm 21.
  \]
\end{enumerate}

\paragraph{四、计算题}

\begin{enumerate}
  \item 令 $y=x^2$，$0\le x\le 1$，则 $dy=2x\,dx$，所以
  \[
    \int_L xy\,dx+(y-x)\,dy
    =
    \int_0^1\left(x^3+(x^2-x)2x\right)\,dx
    =
    \int_0^1(3x^3-2x^2)\,dx
    =
    \frac1{12}.
  \]

  \item 添上顶面 $\Sigma_0:z=1,\ x^2+y^2\le1$ 组成闭曲面。向量场
  \[
    \mathbf F=(x,y,4-2z)
  \]
  的散度为 $1+1-2=0$，故闭曲面积分为 $0$。底面外法向量为 $(0,0,1)$，底面积分为
  \[
    \iint_{\Sigma_0}(4-2)\,dxdy=2\pi.
  \]
  因此锥面下侧部分积分为 $-2\pi$。
\end{enumerate}

\paragraph{五、证明题}

将 $f(x)=x$ 偶延拓到 $[-\pi,\pi]$，其余弦级数为
\[
  x\sim \frac{\pi}{2}
  +\sum_{n=1}^{\infty}\frac{2\big((-1)^n-1\big)}{\pi n^2}\cos nx
  =
  \frac{\pi}{2}
  -\frac{4}{\pi}\sum_{n=1}^{\infty}\frac{\cos(2n-1)x}{(2n-1)^2}.
\]
取 $x=0$ 得
\[
  0=\frac{\pi}{2}
  -\frac{4}{\pi}\sum_{n=1}^{\infty}\frac1{(2n-1)^2},
  \qquad
  \sum_{n=1}^{\infty}\frac1{(2n-1)^2}=\frac{\pi^2}{8}.
\]

\paragraph{六、计算题}

当 $|x|<1$ 时，令 $u=x^2$，有
\[
  \sum_{n=1}^{\infty}nx^{2n}
  =
  \sum_{n=1}^{\infty}nu^n
  =
  \frac{u}{(1-u)^2}
  =
  \frac{x^2}{(1-x^2)^2}.
\]
在 $x=\pm1$ 时级数发散，故收敛域为 $(-1,1)$。

\paragraph{七、证明题}

记
\[
  I=\int_0^2 dx\int_x^2 f(x)f(y)\,dy.
\]
由积分区域关于直线 $x=y$ 对称，交换 $x,y$ 得
\[
  I=\int_0^2 dy\int_0^y f(x)f(y)\,dx.
\]
两式相加，覆盖正方形 $[0,2]\times[0,2]$，因此
\[
  2I=\int_0^2\int_0^2 f(x)f(y)\,dxdy
  =
  \left[\int_0^2 f(x)\,dx\right]^2,
\]
命题得证。

\paragraph{八、证明题}

任取 $N>0$。当 $y\in[-N,N]$ 且 $x\ge1$ 时，
\[
  0\le \frac1{\sqrt[3]{x^4+y^2}}\le \frac1{\sqrt[3]{x^4}}=\frac1{x^{4/3}},
\]
而 $\int_1^{+\infty}x^{-4/3}\,dx$ 收敛。因此由含参广义积分的一致收敛判别，$F(y)$ 在 $[-N,N]$ 上连续。由于 $N$ 任意，$F(y)$ 在 $(-\infty,+\infty)$ 上连续。

\paragraph{九、应用题}

由
\[
  f(x)=\sin x-\int_0^x(x-t)f(t)\,dt
\]
对 $x$ 求导，得
\[
  f'(x)=\cos x-\int_0^x f(t)\,dt.
\]
再求导得
\[
  f''(x)=-\sin x-f(x).
\]
即
\[
  f''+f=-\sin x,\qquad f(0)=0,\quad f'(0)=1.
\]
解得
\[
  f(x)=\frac12\sin x+\frac{x}{2}\cos x.
\]
